\section{Introduction}
\label{sec:introduction}

The conformational changes that matter biologically are collective: a loop
folding over an occupied active site, two domains rotating against one another
so that a duplex groove opens far enough to admit a protein side chain. Tens to
hundreds of atoms move together, over microseconds or longer, and molecular
dynamics reaches that regime slowly for a structural reason rather than an
implementation one. The integrator has to resolve the fastest motion present,
which is the stretch of a bond to hydrogen, so the step is a few femtoseconds
and a microsecond costs of order $10^{9}$ of them \cite{dror2012microscope};
trajectories long enough to watch a domain rearrange or a small protein fold
have generally required purpose-built hardware
\cite{shaw2010anton,lindorfflarsen2011folding}. The usual response is to bias
the dynamics, by umbrella sampling \cite{torrie1977umbrella}, replica exchange
\cite{sugita1999remd}, metadynamics \cite{laio2002metadynamics} or one of their
relatives \cite{bernardi2015enhanced}, which buys sampling at the price of
either many replicas or a collective variable that has to be guessed before the
run starts.

Natural-move Monte Carlo attacks the same problem from the other side. Instead
of displacing every atom by a small amount and repeating, it displaces a
structurally meaningful group of atoms by a large amount, once. The user
partitions the molecule into groups (a chain, a domain, a helical arm, a
backbone segment between two anchors) and nests those groups in a hierarchy. A
move picks a group and applies a rigid rotation and translation to it, or turns
torsions inside it, leaving everything outside where it was. That breaks the
covalent geometry at the two boundaries of the group, and the break is repaired
by a stochastic chain-closure algorithm that finds the alternative closed
geometries, chooses among them and reports the weight needed to keep the
Metropolis test exact \cite{minary2010chainclosure}. The composite proposal is
then accepted or rejected against the Boltzmann factor in the usual way
\cite{sim2012naturalmoves}. The granularity is the point: one accepted move can
shift a whole segment along the path a hinge would take, and the nesting lets
one proposal act at the scale of a domain and the next at the scale of a sugar
pucker. Nothing has to be chosen in advance about \emph{where} the system should
go, only about which parts of it are plausibly rigid. The approach has predicted
sequence-dependent nucleosome occupancy without training on occupancy data
\cite{minary2014nucleosome}, refined structures into single-particle EM and
cryo-EM densities \cite{zhang2012multiscale,chang2020rnacryoem}, modelled
epigenetic marks on DNA \cite{krawczyk2018epigenetic}, and followed peptide
dissociation from MHC class~I \cite{knapp2016pmhc}.

All-atom Monte Carlo on additive force fields is not itself new, and this paper
does not claim it. CAMPARI samples all-atom polypeptide and nucleic-acid models
under the ABSINTH implicit-solvent model and under imported additive parameter
sets \cite{vitalis2009mcmethods}; MCPRO and BOSS were built around OPLS-AA
\cite{jorgensen2005mcpro}; ProtoMS computes binding free energies under AMBER
and OPLS parameters \cite{protoms-software}; and Cassandra samples all-atom
fixed-charge models of molecular fluids \cite{shah2017cassandra}. Concerted
rotation closes a protein or nucleic backbone under OPLS in much the same spirit
as a natural move does
\cite{ulmschneider2003concerted,ulmschneider2004concertedna}. What is new here
is the combination: a nested natural-move proposal, acting at the scale of a
domain on one step and a sugar pucker on the next, driven by a published
additive potential whose evaluation has been checked term by term against the
molecular dynamics codes that distribute it.

Two things have held the natural-move method back, and neither is the sampling.
The first is the energy. Where the natural-move studies above did not use an
experimental restraint they used a knowledge-based or simplified potential, on a
reduced representation where the question allowed one \cite{minary2008foldspace},
the same family of scoring functions Monte Carlo structure prediction has relied
on more generally \cite{leaverfay2011rosetta}. For the questions those studies
asked this was the right choice: a statistical potential is cheap, it tolerates
a coarse representation, and it was trained on the discrimination being
requested, which is whether one model ranks above another. But ranking is all it
returns. A number produced for one system cannot be compared with one produced
for another, or with experiment, and a structure cannot pass between the Monte
Carlo engine and a molecular dynamics run without changing what the energy
means. The second is access. The engine is a command-line program driven by a
plain-text input file, a residue topology file and a parameter deck, and the
move set has to be written down as a hierarchy of atom groups. A PyMOL plugin,
PymoSAICS \cite{pymosaics-plugin}, was written because defining those groups by
hand was the step users found hardest, and it remains a plugin: it lives inside
PyMOL and reaches only the users who already have the engine. Closing both gaps is what this paper
is about.

The energies to adopt are the additive force fields molecular dynamics already
uses, and the specific ones matter, because each was written to repair a named
failure of its predecessor. For protein, ff14SB refits the backbone $\phi/\psi$
and side-chain $\chi$ dihedral terms of ff99SB against quantum reference data
\cite{maier2015ff14sb}; ff19SB then removes the backbone dihedral terms
altogether and replaces them with amino-acid-specific two-dimensional
$(\phi,\psi)$ correction maps fitted to solution-phase surfaces
\cite{tian2020ff19sb}, using the grid-based CMAP cross-term introduced for
CHARMM a decade and a half earlier
\cite{mackerell2004cmap,mackerell2004cmapfull}. For nucleic acids, parmbsc0
refits $\alpha$ and $\gamma$ and so removes the irreversible backbone flips that
used to destroy long DNA trajectories \cite{perez2007bsc0}; parmbsc1 goes
further into $\epsilon/\zeta$, $\chi$ and sugar pucker, benchmarked against a
large set of experimental DNA structures \cite{ivani2016bsc1}. The Olomouc
corrections take a different route: $\chi_{\mathrm{OL3}}$ for RNA
\cite{zgarbova2011ol3}, a $\beta$ refit that gives OL15 for DNA
\cite{zgarbova2015ol15}, an $\alpha/\gamma$ refit against Z-DNA that gives OL21
\cite{zgarbova2021ol21}, and a deoxyribose pucker refit that gives OL24
\cite{zgarbova2025ol24}. The DNA variants among these have been compared
against experiment over long trajectories
\cite{galindomurillo2016amberdna,love2023amberdna}. CHARMM36 is a separately
developed alternative with a different functional form, carrying Urey--Bradley
1--3 terms and harmonic impropers, refit for the DNA BI/BII equilibrium
\cite{hart2012charmm36dna} and for RNA $2'$-hydroxyl sampling
\cite{denning2011charmm36rna}.

These parameter sets are not interchangeable, and they differ from one another
substantially. How much they differ is not something that can be read off an
absolute energy at a single geometry, and this paper does not try: a Fourier
torsion sum carries an arbitrary additive constant, so a refit moves the
constant as freely as it moves the shape, and the difference between two
parameter sets' absolute torsion energies at one conformation is not a physical
quantity. What an all-atom engine buys is narrower and more useful. The number
it returns is the published function's, evaluated on a named parameter set, so a
preference the sampler reports can be attributed to that parameter set, and an
ensemble it produces and an ensemble a molecular dynamics code produces under
the same parameters are ensembles of the same function. A knowledge-based
potential returns a ranking on a scale of its own and can do neither.

This paper reports two things. The first is that MOSAICS evaluates these force
fields at all-atom resolution, and how well. Little of that is new engine code.
MOSAICS reads a force field as a deck of plain-text parameter files, and the
code that evaluates those terms was already there. CHARMM36's Urey--Bradley
1--3 terms go into a slot the bend record already carried, evaluated by the
expression the bend routine already applied to it; its 188 harmonic impropers
go onto the power series the torsion record already expanded. The
correction-map machinery was present and its energy was already correct, but
its analytic force was identically zero at a dihedral of exactly $\pi$, which is
where a peptide built extended sits. That repair, a memory-allocation defect
that ran on every committed job, and a deck reader that carries a per-molecule
label on torsion entries and chains the residues it builds are the changes to
engine code reported here; the released engine cannot read most of these decks
at all. The rest of the work is conversion:
the decks, the converters that write them, and the defects that writing them
exposed. The verification is term by term, on the same coordinates, in vacuum
and with no cutoff, at one fixed conformation per system, against three
independently written programs, sander~24.0, OpenMM~8.5.2 on the Reference
platform and GROMACS~2025.4 in double precision
\cite{case2023ambertools,eastman2017openmm7,eastman2024openmm8,abraham2015gromacs},
following the cross-engine practice established for this purpose
\cite{shirts2017engines,braun2018bestpractices}. It covers eleven systems of
56 to 6,139 atoms and 26 system-and-force-field combinations
(Figure~\ref{fig:systems}), and it is completed by a finite-difference check of
the analytic forces and by a measurement of what one Monte Carlo step costs.

The second is \studio{}, which puts that engine in front of a user who has
never seen its input syntax, in a web browser. The engine is compiled to
WebAssembly and the Studio's own pipeline, four operations, sequence to
structure, structure to minimised structure, variant building and structure
comparison, runs beside it under Pyodide, so a sequence or a coordinate file
goes in, a minimised structure and a molecular graph come out, nothing is
installed and nothing leaves the visitor's machine. Every run carries a record
naming the engine build, the parameter files and the inputs by digest. The
WebAssembly engine was checked against the native engine on a documented
minimisation, and that number is reported alongside the force-field ones,
because the Studio is only as good as the engine it carries.

After every electrostatic term is put on a single Coulomb constant, nineteen of
the 26 combinations agree with every reference available to them to within
0.001~\kcal{}, the median largest gap over the 26 is 0.00037~\kcal{}, and the
one combination that is worse than 0.01~\kcal{} is BPTI under ff19SB, where a
single improper is read in a different atom order by the AMBER topology and the
MOSAICS residue template. On the two systems compared term by term, the largest
difference between MOSAICS and OpenMM on any term is $1.8\times10^{-4}$~\kcal{}
on the 818-atom DNA:RNA heteroduplex and $9.6\times10^{-6}$~\kcal{} on the
capped peptide, in both cases on the angle term, and in both cases traceable to
the reference topology and not to MOSAICS. Analytic forces were checked against
a central difference of the same energy in 61 term-and-system combinations;
fitted over a fixed window, the convergence exponent lies between 1.994 and
2.009. One Monte Carlo step at 818 atoms costs 3.46~ms on one core of a laptop.

Section~\ref{sec:methods} describes the systems, the parameter decks, how each
reference calculation was produced, and how the Studio is built and served.
Section~\ref{sec:results} gives the 26-cell comparison, the term-by-term
comparison for the two reference systems, the electrostatic-constant correction,
the force checks, the cost measurement and the Studio's own verification.
Section~\ref{sec:discussion} sets out what follows from single-point agreement,
and what does not.
